% Figura 6 del consolidado académico (PDF vigente, página 6): cobertura de
% etapas, intersecciones funcionales y señales positivas de evidencia.
% Reconstruida en LaTeX a partir de los valores impresos en el original.
\begin{figure}[!htbp]
  \centering
  \begin{minipage}[t]{0.47\linewidth}
    \centering
    {\sffamily\bfseries\fontsize{7.6}{8.2}\selectfont (a) SP1 concentra la cobertura\par}
    \vspace{1.5mm}
    \resizebox{\linewidth}{!}{%
\begin{tikzpicture}[x=1mm,y=1mm,font=\sffamily\fontsize{6.0}{6.5}\selectfont]
% La escala deja sitio al rótulo del valor más alto dentro del cuadro:
% 115 documentos ocupan 55 mm y la etiqueta cabe antes del borde en 100 mm.
\path[use as bounding box] (0,-9) rectangle (100,34);
\def\sc{0.478}
\foreach \y/\lab/\v/\col in {26/{SP1 · coalición / asignación}/115/{gamecol},14/{SP3 · planificación / tráfico}/20/{canary!80!onyx},2/{SP2 · transporte físico}/15/{autumn}}{%
  \fill[onyx!5] (36,\y-3.4) rectangle (36+\sc*122,\y+3.4);
  \fill[\col,rounded corners=.5pt] (36,\y-3.4) rectangle (36+\sc*\v,\y+3.4);
  \node[anchor=east,text=ink] at (34.5,\y) {\lab};
  \node[anchor=west,font=\bfseries,text=ink] at (36+\sc*\v+1.4,\y) {\v};}
\draw[-{Stealth[length=1.4mm]},ink,line width=.35pt] (36,-3)--(36+\sc*124,-3);
\node[anchor=north,text=muted,font=\fontsize{5.2}{5.6}\selectfont] at (36+\sc*62,-4){documentos del corpus};
\end{tikzpicture}}
  \end{minipage}\hfill
  \begin{minipage}[t]{0.50\linewidth}
    \centering
    {\sffamily\bfseries\fontsize{7.6}{8.2}\selectfont (b) La integración entre etapas es excepcional\par}
    \vspace{1.5mm}
    \resizebox{\linewidth}{!}{%
\begin{tikzpicture}[x=1mm,y=1mm,font=\sffamily\fontsize{5.8}{6.2}\selectfont]
% La escala vertical deja 6 mm libres sobre la barra de 103 para su rótulo y
% para la nota, de modo que el panel no invade al de la izquierda.
\path[use as bounding box] (-6,-14) rectangle (104,40);
\def\sy{0.245}
\foreach \i/\lab/\v in {0/{solo\\SP1}/103,1/{solo\\SP2}/13,2/{solo\\SP3}/10,3/{SP1 +\\SP3}/10,4/{SP1 +\\SP2}/2,5/{SP2 +\\SP3}/0,6/{las\\3}/0}{%
  \pgfmathsetmacro{\x}{10+\i*13}
  \pgfmathsetmacro{\h}{\sy*\v}
  \ifdim \v pt>0pt
    \fill[onyx!42,rounded corners=.5pt] (\x,0) rectangle (\x+9,\h);
  \fi
  \node[anchor=south,font=\bfseries\fontsize{5.8}{6.2}\selectfont,text=ink] at (\x+4.5,\h+0.8) {\v};
  \node[anchor=north,align=center,text=muted,font=\fontsize{5.0}{5.4}\selectfont] at (\x+4.5,-6.2) {\lab};}
% marcadores de etapa bajo cada columna
\foreach \i/\a/\b/\c in {0/1/0/0,1/0/1/0,2/0/0/1,3/1/0/1,4/1/1/0,5/0/1/1,6/1/1/1}{%
  \pgfmathsetmacro{\x}{10+\i*13}
  \foreach \k/\col/\on in {0/{gamecol}/\a,1/{autumn}/\b,2/{canary!80!onyx}/\c}{%
    \pgfmathsetmacro{\yy}{-1.6-\k*1.9}
    \ifnum\on=1 \fill[\col] (\x+4.5,\yy) circle (0.75);
    \else \draw[onyx!35,line width=.22pt] (\x+4.5,\yy) circle (0.75);\fi}}
\node[anchor=east,text=muted,font=\fontsize{4.6}{5.0}\selectfont] at (8,-1.6){SP1};
\node[anchor=east,text=muted,font=\fontsize{4.6}{5.0}\selectfont] at (8,-3.5){SP2};
\node[anchor=east,text=muted,font=\fontsize{4.6}{5.0}\selectfont] at (8,-5.4){SP3};
\node[anchor=north west,align=left,text width=62mm,text=muted,font=\fontsize{5.0}{5.6}\selectfont] at (30,38)
  {103 trabajos cubren solo SP1; solo 2 enlazan SP1--SP2 y ninguno cubre las tres etapas.};
\node[rotate=90,anchor=south,text=muted,font=\fontsize{5.2}{5.6}\selectfont] at (-4,14){documentos};
\end{tikzpicture}}
  \end{minipage}

  \vspace{2.5mm}
  {\sffamily\bfseries\fontsize{7.6}{8.2}\selectfont (c) La evidencia física y formal aparece con mucha menor frecuencia\par}
  \vspace{1.5mm}
  \makebox[\textwidth][c]{\resizebox{0.86\textwidth}{!}{%
\begin{tikzpicture}[x=1mm,y=1mm,font=\sffamily\fontsize{6.0}{6.5}\selectfont]
\def\sc{1.30}
\foreach \y/\lab/\v/\col in {24/{Simulación}/99/{learncol},16/{Contexto / caso industrial}/25/{mustard!85!onyx},8/{Experimento físico}/12/{autumn},0/{Análisis formal}/7/{onyx!45}}{%
  \fill[onyx!5] (40,\y-2.6) rectangle (40+\sc*105,\y+2.6);
  \fill[\col,rounded corners=.5pt] (40,\y-2.6) rectangle (40+\sc*\v,\y+2.6);
  \node[anchor=east,text=ink] at (38.5,\y) {\lab};
  \node[anchor=west,font=\bfseries,text=ink] at (40+\sc*\v+1.4,\y) {\v};}
\foreach \v in {0,20,40,60,80,100}{%
  \node[anchor=north,text=muted,font=\fontsize{5.2}{5.6}\selectfont] at (40+\sc*\v,-3.4){\v};}
\node[anchor=north west,text=muted,font=\fontsize{5.0}{5.4}\selectfont] at (40,-7)
  {«Contexto / caso industrial» puede acompañar a simulación; son conteos positivos, no prevalencias.};
\end{tikzpicture}%
  }}
  \caption[Cobertura de etapas, intersecciones y señales de evidencia.]{Cobertura
  de etapas, intersecciones funcionales y señales positivas de evidencia en el
  corpus analítico. (a) Cobertura codificada por etapa. (b) Intersecciones
  observadas bajo la taxonomía SP1--SP3; los ceros significan «no identificado
  bajo esta codificación», no inexistencia bibliográfica. (c) Conteos positivos
  por tipo de evidencia; no se presentan porcentajes de prevalencia porque la
  disponibilidad de extracción estructural varía por campo.}
  \label{fig:lit-stage-coverage}
  \viusource{tablas derivadas del mapeo analítico, $n=244$}
\end{figure}
